Artificial subsurface (tile) drainage is among the most consequential
anthropogenic modifications of the field-to-stream pathway, draining roughly
$12$--$17\%$ of the world's cultivated land and more than $22$~Mha in the
United States alone \parencite{Valipour2014,Valayamkunnath2020}. This review consolidates the state of knowledge on
drained-land hydrology, water quality, agronomy, and management, emphasising
well-supported findings, mechanisms, and quantified magnitudes over open
questions. Its unifying conclusion is that drainage is fundamentally a
hydrological intervention: by intercepting water that would otherwise recharge
groundwater or leave as surface runoff and routing it rapidly to streams, tiles
commonly carry $40$--$65\%$ of streamflow in humid settings, dominate field-scale
water export, elevate low flows and baseflow, shorten travel times, and
homogenise catchment response \parencite{King2014,Schilling2019}. This hydrology governs the water-quality
outcomes that follow. Relative to undrained or surface-drained land, subsurface
drainage reduces runoff, sediment, and phosphorus but opens an efficient conduit
for nitrate, the dominant and best-quantified export; nitrate loads are
hydrologically controlled, strongly seasonal, and concentrated in high-flow
episodes, with concentrations routinely near or above the $10~\mathrm{mg~L^{-1}}$
drinking-water limit. The presumption of a filtering ``clean'' subsurface pathway
has been overturned: a small, hydraulically active volume of macropores and
cracks delivers particulate phosphorus and sediment to drains within hours \parencite{Pease2018}, so
tiles frequently dominate edge-of-field phosphorus loads \parencite{Pease2018,Sharpley1995}. Against these losses,
a tiered hierarchy of in-field, drainage-design, and edge-of-field practices
produces quantified but scale-attenuated reductions---cover crops clustering
around $40$--$60\%$ nitrate-load reduction \parencite{Malone2018,Malone2014b,Malone2023}, alongside controlled drainage,
bioreactors, wetlands, and saturated buffers governed by residence time,
temperature, and loading \parencite{Christianson2013,Jaynes2008}---so meeting downstream targets requires stacking
complementary measures. Drainage also reliably raises yields, reclaims saline
and waterlogged land, and returns favourable economic and social benefits, while
validated process models, geophysics, and remote sensing underpin design and
mapping. Because independent methods, climates, and scales converge on the same
mechanisms and comparable magnitudes, these conclusions can be stated with
confidence, and recognising drainage water quality as a design criterion is
their consolidated implication.